The polynomials $P_n(\rho_1,\ldots,\rho_n)$ are not uniquely determined by equation~\eqref{equ.polinomigeneratori}. In fact, according to formula~\eqref{equ.bracketmuenne}, we have
\begin{gather}\label{equ.bracketrhoenne}
[\rho_n,\rho_m]=(n-m)\frac{(n+m+1)!}{(n+1)!(m+1)!}\rho_{n+m},\end{gather}
and therefore, by Poincar\'e--Birkhof\/f--Witt theorem, we can choose every $P_n$ as a linear combination with rational coef\/f\/icients of monomials of type
\begin{gather*} \rho_1^{i_1}\rho_2^{i_2}\cdots \rho_n^{i_n},\qquad \text{with}\qquad\sum_{h=1}^n hi_h=n .\end{gather*}

The following theorem, which is mathematically nontrivial, tells us in particular that, in the (super) commutative case, we can restrict to monomials such that $i_2+\cdots+i_n\le 1$.

\begin{Theorem}\label{thm.standardform}
In the notation above, for every integer $n>0$, there exists an unique sequence
$c_1^n,\ldots,c_n^n$ of rational numbers
such that, for every linear operator $f\colon A\to A$, we have
\begin{gather*}\Phi^{n+1}_f=\big(c_1^n\rho_1^n+c^n_2\rho_1^{n-2}\rho_2+c^n_3\rho_1^{n-3}\rho_3+\cdots+c^n_n\rho_n\big)f.\end{gather*}
\end{Theorem}
